\documentclass[12pt,a4paper]{report}

\usepackage{graphicx, subfigure}
\usepackage{appendix}

\newcommand{\HRule}{\rule{\linewidth}{0.5mm}}

\begin{document}
\begin{titlepage}

\begin{center}

% Upper part of the page
\includegraphics[scale=0.4]{figures/logo}\\[1cm]

\textsc{\LARGE Advanced Networking 2011-2012}\\[1.5cm]

\textsc{\Large Final Report}\\[0.5cm]

% Title
\HRule \\[0.5cm]
{\huge \bfseries Peer Streamer yada yada...}\\[0.5cm]
%--------------------------------------------------------------------------\\[0.4cm]
%{\LARGE \bfseries Trusted Organizational Base}\\
\HRule \\[1.5cm]

% Author and supervisor
\begin{minipage}{0.4\textwidth}
\begin{flushleft} \large
Arb\"er \textsc{Fama}\\
\emph{arberfama@gmail.com}
\end{flushleft}
\end{minipage}
\begin{minipage}{0.4\textwidth}
\begin{flushright} \large
Giovanni \textsc{Simoni} \\
\emph{simgidacav@gmail.com}
\end{flushright}
\end{minipage}

\vfill

% Bottom of the page
{\large \today}

\end{center}

\end{titlepage}


\begin{abstract}
In this report it is explained a module extension for \textbf{GRAPES} library.
It consists of a practical implementation of a module that manages the
transmittion over TCP. In the previous version of GRAPES the transmission was
executed over UDP. On the one hand, the advantages of TCP over UDP are quite
clear, as TCP guarantees that the sent data actually arrives, that it arrives
in order and that there are no duplicates, while UDP provides none of these
guarantees (only offers best effort policy). Unlike TCP, UDP does not provide
any flow and congestion control. On the other hand, UDP is more suitable in
applications and enviroments such peer-to-peer streaming, where thorughput
performance is the main request. It is faster because there is no
error-checking for packets, it does not order packets and there is no
tracking connections. It is a small transport layer designed on top of IP.
If any ordering is required, it has to be managed by the application layer.\\
Considering these aspects, there have been choose to develop a module that
consideres advantages from both protocols and takles some drawbacks derived
from them.

\end{abstract}

\tableofcontents
%\listoffigures
%\listoftables

\chapter{Introduction}
\label{ch:intro}
Before getting lost in many technicalities, there have been reported a small
introduction to the GRAPES library. It is a generic environment for P2P
streaming. It is designed to be usable in different environments and
situations. Peer-to-peer (\textbf{P2P}) tecnologies are becoming increasingly
popular as a way to overcome the scalability problem limitations intrinsic in
traditional applications based on client/server paradigm. In particular there
is a great interest in P2P streaming applications because they have high
demands in term of bandwidth requirements and IP-level multicasting is not
supported on the internet.
More details could be found in \cite{disi10-038}.

In this report ....
In particular, we will focus on:
\begin{itemize}
  \item yada
  \item yada
  \item yada
\end{itemize}
More details about \textbf{yada} can be found in \cite{yada1}, \cite{yada2}, \cite{yada3} and
\cite{yada4}.
The first approach for studying the structure of these entities is by using yada...
That is:
\begin{itemize}
\item a
\item b
\item c
\end{itemize}

%=========================================================================================================================

\chapter{GRAPES}
\label{ch:grapes}
At this point lets have a look to the content of \textbf{GRAPES} library and the module we developed.

\section{asd}
\label{asd}

%=========================================================================================================================

\chapter{Tests}
\label{ch:tests}
In this chapter we are going to describe the tests ....

%=========================================================================================================================

\chapter{Performances Analysis}
In this chapter we are going to annalyse the performances...

\begin{center}
    \begin{tabular}{ | p{4.8cm} | p{4.8cm} | p{4.8cm} |}
    \hline
       \textbf{column 1}
         & \textbf{column 2}
         & \textbf{column 3}\\ \hline \hline
       .....
         & ------
         & .... \\ \hline
     \end{tabular}
\end{center}

%=========================================================================================================================
\chapter{Conclusion}

%=========================================================================================================================

\begin{appendices}

\section{Net-helper TCP}

\subsection{Send to peer}
\label{nh-send2peer}
\begin{lstlisting}
int send_to_peer(const struct nodeID *self, struct nodeID *to,
                 const uint8_t *buffer_ptr, int buffer_size)
{
    client_t cl;
    local_t *loc;
    struct timeval zero = {
        .tv_sec = 0,
        .tv_usec = 0
    };
    const msgbuf_t msg = {
        .data = buffer_ptr,
        .size = buffer_size
    };

    assert(self->loc != NULL);

    loc = self->loc;

    cl = dict_search(loc->neighbors, &to->addr);
    if (!client_valid(cl)) {
        int clfd;

        if (client_connect(cl, &to->addr) == -1) {
            return -1;
        }
        clfd = client_get_fd(cl);
        if (sockaddr_send_hello(&self->addr, clfd) == -1) {
            close(clfd);
            return -1;
        }
    }

    if (client_write(cl, &msg) == -1) return -1;
    if (aqueue_scan_dict(loc->aqueue, loc->neighbors,
                         loc->srvfd, NULL, &zero) == -1) {
        return -1;
    }
    return buffer_size;
}
\end{lstlisting}

\subsection{Receive from peer}
\label{nh-rcv-peer}
\begin{lstlisting}
int recv_from_peer(const struct nodeID *self, struct nodeID **remote,
                   uint8_t *buffer_ptr, int buffer_size)
{
    local_t *loc;
    client_t sender;
    nodeid_t *sender_node;
    const msgbuf_t *msg;

    assert(self->loc != NULL);
    loc = self->loc;

    while (aqueue_empty(loc->aqueue)) {
        if (aqueue_scan_dict(loc->aqueue, loc->neighbors, loc->srvfd,
                             NULL, NULL) == -1) {
            return -1;
        }
    }

    sender = aqueue_next(loc->aqueue);
    msg = client_read(sender);
    assert(msg != NULL);

    if (msg->size > buffer_size) {
        print_err(``recv_from_peer'', NULL, ENOBUFS);
        return -1;
    }
    memcpy((void *)buffer_ptr, msg->data, msg->size);

    sender_node = create_node(NULL, 0);
    sockaddr_copy(&sender_node->addr, client_get_remote(sender));
    *remote = sender_node;

    return (int) msg->size;
}
\end{lstlisting}

\subsection{Wait for data}
\label{nh-wait4data}
\begin{lstlisting}
int wait4data(const struct nodeID *self, struct timeval *tout,
              int *user_fds)
{
    local_t *loc;
    struct timeval tmp;
    int res, ret;
    tout_t T;

    assert(self->loc != NULL);
    loc = self->loc;

    T = tout_new(tout);
    ret = 0;
    while (!tout_expired(T) && aqueue_empty(loc->aqueue)) {
        res = aqueue_scan_dict(loc->aqueue, loc->neighbors, loc->srvfd,
                               NULL, tout_remaining(T, &tmp));
        if (res == -1) {
            tout_del(T);
            return -1;
        }
        ret |= res;
    }
    tout_del(T);

    return ret || !aqueue_empty(loc->aqueue);
}
\end{lstlisting}


\subsection{Other functions}
\label{nh-oth-func}
\begin{lstlisting}
int tcp_serve (int backlog, sockaddr_t *addr)
{
    int fd;

    fd = socket(AF_INET, SOCK_STREAM, 0);
    if (fd == -1) {
      print_err(``tcp_serve'', ``socket'', errno);
        return -1;
    }

    if (bind(fd, &addr->sa, sizeof(struct sockaddr_in))) {
        print_err(``tcp_serve'', ``bind'', errno);
        close(fd);
        return -1;
    }

    if (listen(fd, backlog) == -1) {
        print_err(``tcp_serve'', ``listen'', errno);
        close(fd);
        return -1;
    }

    return fd;
}
\end{lstlisting}

\begin{lstlisting}
struct nodeID *nodeid_dup (struct nodeID *s)
{
    s->refcount ++;
    return s;
}
\end{lstlisting}

\begin{lstlisting}
int nodeid_equal (const nodeid_t *s1, const nodeid_t *s2)
{
    if (s1 == s2) return 1;
    return sockaddr_equal(&s1->addr, &s2->addr);
}
\end{lstlisting}

\begin{lstlisting}
struct nodeID *create_node (const char *ipaddr, int port)
{
    nodeid_t *ret = mem_new(sizeof(struct nodeID));

    /* Initialization of address */
    memset(&ret->addr, 0, sizeof(sockaddr_t));
    ret->addr.sa.sa_family = AF_INET;
    sockaddr_set_port(&ret->addr, port);

    if (ipaddr != NULL) {
        if (sockaddr_in_init(&ret->addr, ipaddr, port) == -1) {
            return NULL;
        }
    }

    ret->loc = NULL;
    ret->refcount = 1;

    return ret;
}
\end{lstlisting}

\begin{lstlisting}
void nodeid_free (struct nodeID *s)
{
    if (s == NULL) return;

    s->refcount --;
    if (s->refcount == 0) {
        local_t *local = s->loc;

        if (local != NULL) {
            dict_del(local->neighbors);
            if (local->srvfd != -1) close(local->srvfd);
            aqueue_del(local->aqueue);
            free(local);
        }
        free(s);
    }
}
\end{lstlisting}

\begin{lstlisting}
struct nodeID * net_helper_init (const char *ipaddr, int port,
                                 const char *config)
{
    nodeid_t *self;
    struct tag *cfg_tags;
    int backlog;
    local_t *local;

    self = create_node(ipaddr, port);
    if (self == NULL) {
        return NULL;
    }

    self->loc = local = mem_new(sizeof(local_t));

    /* Default settings */
    backlog = DEFAULT_BACKLOG;

    /* Reading settings */
    cfg_tags = NULL;
    if (config) {
      cfg_tags = config_parse(config);
    }
    if (cfg_tags) {
        /* FIXME: this seems not to work! Testing needed */
        config_value_int_default(cfg_tags, CONF_KEY_BACKLOG, &backlog,
                                 DEFAULT_BACKLOG);
    }
    local->neighbors = dict_new(cfg_tags);
    free(cfg_tags);

    local->aqueue = aqueue_new();
    local->srvfd = tcp_serve(backlog, &self->addr);
    if (local->srvfd == -1) {
      nodeid_free(self);
        return NULL;
    }

    return self;
}
\end{lstlisting}

\begin{lstlisting}
struct nodeID *nodeid_undump (const uint8_t *b, int *len)
{
    nodeid_t *node;

    node = create_node(NULL, 0);
    *len = sockaddr_undump(&node->addr, sizeof(sockaddr_t),
                           (const void *)b);
    return node;
}
\end{lstlisting}

\begin{lstlisting}
int nodeid_dump (uint8_t *b, const struct nodeID *s,
                 size_t max_write_size)
{
    return sockaddr_dump((void *)b, max_write_size, &s->addr);
}
\end{lstlisting}

\begin{lstlisting}
int node_ip(const struct nodeID *s, char *ip, int len)
{
    return sockaddr_strrep(&s->addr, ip, len) == NULL ? 0 : -1;
}
\end{lstlisting}

\begin{lstlisting}
int node_ip(const struct nodeID *s, char *ip, int len)
{
    return sockaddr_strrep(&s->addr, ip, len) == NULL ? 0 : -1;
}
\end{lstlisting}


\end{appendices}

\bibliographystyle{plain}
\begin{thebibliography}{99}

\bibitem{disi10-038} \textit{Luca Abeni, Csaba Kiraly, Alessandro Russo,
Marco Biazzini and Renato Lo Cigno},
June 2010, \textbf{GRAPES: a Generic Environment for P2P Streaming.}
{\em Technical report \#DISI-10-38}

\end{thebibliography}
%\bibliography{rep-bib}

\end{document}
